\FifthArticle{20}
For further illustration, the calculation of the imaginary roots of the equation treated above,
\[
x^7+28x^4-480=0
\]
may serve as an example. In the notation of Article 17 we first have, as above, $m=4,\ n=3,\ e=28,\ f=480$, and further $\varepsilon=0,\ \varphi=180^{0}$. The formulas I of Article 18 therefore become
\[
\begin{aligned}
r^4&=\frac{480\sin7\rho}{28\sin3\rho}\\
r^7&=-\frac{480\sin4\rho}{\sin3\rho}\\
r^3&=-\frac{28\sin4\rho}{\sin7\rho}
\end{aligned}
\]
and formula II becomes
\[
\frac{1}{\lambda}=\frac{823543}{6750}=\frac{\sin7\rho^7}{\sin3\rho^3\sin4\rho^4}
\]
From this equation two values lying between $0$ and $90^{0}$ must be determined for $\rho$, since the equation $\mathrm{X}=0$, besides its three real roots already found, has two pairs of imaginary roots belonging together. Within these limits, $\sin7\rho$ $=0$ at three values, namely $\rho=25\frac57$ degrees, $51\frac37$ degrees and $77\frac17$ degrees; $\sin7\rho^7$ changes sign each time. Likewise, $\sin3\rho$ $=0$ once, for $\rho=60^{0}$, with a change of sign in $\sin3\rho^3$. Finally, $\sin4\rho$ $=0$ once, for $\rho=45^{0}$, but without a change of sign in $\sin4\rho^4$. Taking into account also that the value of $\frac{\sin7\rho^7}{\sin3\rho^3\sin4\rho^4}$ for $\rho=0$ is to be set equal to the limiting value $\frac{7^7}{3^3 4^4}$, the behaviour of this function in the six subdivisions of the interval from $0$ to $90^{0}$ is summarized in the following table:
\SourceLeaf{107}
\[
\begin{array}{c|c}
\rho=0\ \text{degrees}&\frac{823543}{6912}\\
&+\\
25\frac57&0\\
&-\\
45&\infty\\
&-\\
51\frac37&0\\
&+\\
60&\infty\,.\\
&-\\
77\frac17&0\\
&+\\
90^{0}&\infty
\end{array}
\]
It follows that a value of $\rho$ satisfying formula II must lie in each of the fourth and sixth intervals. Nothing further is needed for our purpose, since it is already known from the outset that there are only two such values. I put equation II in the form
\[
7\log\sin7\rho-3\log\sin3\rho-4\log\sin4\rho=S=2{,}0863825
\]
The following scheme shows how the value lying between $51\frac37$ and $60$ degrees is found by successive approximation with tables of $3,\ 4,\ 5,\ 7$ digits:
\[
\begin{array}{c|c|c}
\rho&S&\text{Error}\\
57^{0}&1{,}527&-0{,}559\\
58&2{,}354&+0{,}268
\end{array}
\]
\[
\begin{array}{c|c|c}
57^{0}40'&2{,}0624&-0{,}0240\\
57\ 50&2{,}2057&+0{,}1193
\end{array}
\]
\[
\begin{array}{c|c|c}
57^{0}41'&2{,}07658&-0{,}00980\\
57\ 42&2{,}09074&+0{,}00436
\end{array}
\]
\[
\begin{array}{c|c|c}
57^{0}41'41''&2{,}0862962&-0{,}0000863\\
57\ 41\ 42&2{,}0865320&+0{,}0001495
\end{array}
\]
This gives $\rho=57^{0}41'41''366$, and then, from the second of the formulas I,
\SourceLeaf{108}
\[
\begin{array}{lrcl}
\log\sin4\rho&&=&9{,}8891425\,n\\
\mathrm{Compl.}\,\log\sin3\rho&&=&0{,}9193523\\
\log(-480)&&=&2{,}6812412\,n\\\hline
7\log r&&=&3{,}4897360\\
\log r&&=&0{,}4985337
\end{array}
\]
and hence
\[
x=+1{,}6843159+2{,}6637914\,i
\]
as well as the other root belonging to it,
\[
x=+1{,}6843159-2{,}6637914\,i
\]

The other value lying between $77\frac{1}{7}$ and $90^{0}$ for $\rho$ is found, by using tables of three decimal places, to lie between $86^{0}$ and $87^{0}$. The calculation, in the same form as in the preceding case, is:
\[
\begin{array}{c|r|r}
\rho&S&\mathrm{Error}\\
86^{0}&1{,}885&-0{,}201\\
87&2{,}533&+0{,}447\\
\end{array}
\]
\[
\begin{array}{c|r|r}
86^{0}10'&1{,}9907&-0{,}0957\\
86\;20&2{,}0946&+0{,}0082\\
\end{array}
\]
\[
\begin{array}{c|r|r}
86\;19&2{,}08409&-0{,}00229\\
86\;20&2{,}09447&+0{,}00809\\
\end{array}
\]
\[
\begin{array}{c|r|r}
86^{0}19'13''&2{,}0863229&-0{,}0000596\\
86\;19\;14&2{,}0864970&+0{,}0001145
\end{array}
\]
\[
\rho=86^{0}19'13''342
\]
\[
\begin{array}{lrcl}
\log\sin4\rho&&=&9{,}4049540\,n\\
\mathrm{Compl.}\,\log\sin3\rho&&=&0{,}0081108\,n\\
\log(-480)&&=&2{,}6812412\,n\\\hline
7\log r&&=&2{,}0943060\,n\\
\log r&&=&0{,}2991866\,n
\end{array}
\]
If one prefers $r$ to be positive, it is only necessary at the same time to give $\rho$ the value increased by $180^{0}$, namely $266^{0}19'13''342$. The root itself is
\[
x=-0{,}1278113-1{,}9874234\,i
\]
and the other root belonging to it differs only in the sign of its imaginary part.
\par\SourceLeaf{109}
All the roots of the equation $x^7+28x^4-480=0$ are therefore
\[
\begin{array}{l}
+1{,}9228841\\
-2{,}4580892\\
-2{,}5778036\\
+1{,}6843159+2{,}6637914\,i\\
+1{,}6843159-2{,}6637914\,i\\
-0{,}1278113+1{,}9874234\,i\\
-0{,}1278113-1{,}9874234\,i
\end{array}
\]
The sum of the roots, $+0{,}0000005$, agrees with the true value $0$ as accurately as could be expected from the use of logarithms of seven digits. In the other form one has
\[
\begin{array}{r|r}
\log r&\rho\\\hline
0{,}2839531&0\\
0{,}3905976&180^{0}\\
0{,}4112498&180\\
0{,}4985337&57\;41'41''366\\
0{,}4985337&302\;18\;18{,}634\\
0{,}2991866&93\;40\;46{,}658\\
0{,}2991866&266\;19\;13{,}342
\end{array}
\]
The sum of the logarithms of the values of $r$ is found to be $=2{,}6812411$, likewise agreeing with satisfactory accuracy with the logarithm of $480$.

It is hardly necessary to recall that the calculations set out in this article and in Article 16 are intended only to illustrate the main stages of the work, and are by no means to prescribe the form of the minor mechanics of the operations. Those with more experience in calculation will usually prefer fewer intermediate stages than were used in those examples. In general, anyone with some experience in such work will readily arrange the details in the form best suited to the circumstances of each case and to their own individual habits. This is not the place to enter further into such details.
\begin{center}\rule{0.25\linewidth}{0.3pt}\end{center}
